% !TeX spellcheck = en_US
\documentclass[french]{yLectureNote}

\title{Outils Mathématiques 2}
\subtitle{Physique}
\author{Paulhenry Saux}
\date{\today}
\yLanguage{Français}

\professor{F.Pettinari}
\usepackage{graphicx}%----pour mettre des images
\usepackage[utf8]{inputenc}%---encodage
\usepackage{geometry}%---pour modifier les tailles et mettre a4paper
%\usepackage{awesomebox}%---pour les boites d'exercices, de pbq et de croquis ---d\'esactiv\'e pour les TP de PC
\usepackage{tikz}%---pour deiffner + d\'ependance de chemfig
% \usepackage{tabularx}%---pour dimensionner automatiquement les tableaux avec variable X
\usepackage{awesomebox}%---Pour les boites info, danger et autres
\usepackage{menukeys}%---Pour deiffner les touches de Calculatrice
\usepackage{fancyhdr}%---pour les en-t\^ete personnalis\'ees
\usepackage{blindtext}%---pour les liens
\usepackage{hyperref}%---pour les liens (\`a mettre en dernier)
\usepackage{caption}%---pour la francisation de la l\'egende table vers Tableau
\usepackage{pifont}
\usepackage{array}%---pour les tableaux
\usepackage{lipsum}
\usepackage{yFlatTable}
\usepackage{multicol}
\newcommand{\Lim}[1]{\lim\limits_{\substack{#1}}\:}
\renewcommand{\vec}{\overrightarrow}
\newcommand{\N}[0]{\mathbb{N}}
\newcommand{\dd}{\mathrm{d}}
\newcommand{\norm}[1]{||\vec{#1}||}
\newcommand{\fo}{\psi(\vec{r},t)}
\newcommand{\foe}{\psi(\vec{r},t)\*}
\newcommand{\HH}{\hat{H}}
\newcommand{\hb}{\hbar}
\newcommand{\lap}{\nabla^2}
\newcommand{\lapcc}{\frac{\partial^2 }{\partial x^2}+\frac{\partial^2 }{\partial y^2}+\frac{\partial^2 }{\partial z^2}}
\newcommand{\mpsi}{\(\psi\)}
\begin{document}
\chapter{Équations différentielles}
\section{Premier ordre}
On cherche la solution générale de l'équation \[y'(x)-\alpha(x)y(x) = f(x)\]. La solution est la somme de la solution de l'équation homogène associée et la solution particulière de cette équation.
\subsection{Solution de l'équation homogène associée}
\begin{theorem}[Solution d'une EQD du premier ordre]
La solution d'une équation de la forme \(y'(x)-\alpha(x)y(x) = 0\) est \[Ce^{\int^x \alpha(t)\dd t}\]
\end{theorem}
\subsection{Méthode de la variation de la constante}
\begin{theorem}[Solution particulière d'une EQD du premier ordre inhomogène]
C'est \[Ce^{\int^x \alpha(t)\dd t}\] avec \[C(x) = \int^x f(y)e^{-\int^y \alpha(t)\dd t}\dd y\]
\end{theorem}
\checkInfo{En pratique}{On calcule d'abord \(h(y) = -\int^y \alpha(t)\dd t\) puis on injecte le résultat dans la solution \(\int^x f(y)e^{h(y)}\dd y\)}

\section{Second ordre}
On cherche à résoudre une équation de la forme \[y''(x)-\alpha(x)y'(x)-\beta(x)y(x) = f(x)\]
\subsection{Solution de l'équation homogène associée}
\subsubsection{Première solution}
On cherche une première sokution ``évidente'', dont la forme sera donnée. Cette solution est notée \(y_1(x)\).
\subsubsection{Seconde solution}
On a une équation de la forme \[y_1(x)y_2'(x)-y_2(x)y_1'(x) = C e^{\int^x \alpha(t)\dd t}\]
avec \(y_1(x), y_1'(x)\) qui sont connus et jouent le r\^ole des coefficients dans cette équation premier ordre que l'on peut résoudre par la section précédente.

\warningInfo{Remarque}{On peut aussi trouver \(y_2(x)\) en calculant \[y_2(x) = y_1(x)\int^x \frac{C e^{\int^t \alpha(m)\dd m}}{y_1^2}\dd t \]}
\subsubsection{Conclusion}
La solution générale est ainsi \[c_1y_1(x)+c_2y_2(x) = y(x)\] avec des constantes à déterminer avec les conditions initiales.
\subsubsection{Solution particulière}
Une solution particulière est de la forme \[y_p(x) = c_1(x)y_1(x)+c_2(x)y_2(x)\] avec
\begin{theorem}[Forme des coeffcients de l'EQD inhomogène du second ordre]
\[c_1(x) = -\int^x \frac{f(t)y_2(t)}{C e^{\int^t \alpha(m)\dd m}}\]
\[c_2(x) = \int^x \frac{f(t)y_1(t)}{C e^{\int^t \alpha(m)\dd m}}\]
\end{theorem}
\subsection{Conclusion}
La solution de l'équation est donc \(c_1y_1(x)+c_2y_2(x)+y_p(x)\)
\tipsInfo{Remarque}{L'expression \(C e^{\int^t \alpha(m)\dd m}\) peut \^etre notée \(W(x)\) et se nomme le Wronskien. On le calcule en premier au début de la résolution pour s'en servir dans les étapes suivantes de la résolution.}
 \end{document}
